\begin{itemize}
\item What it is: a rule saying which subsets $S \subseteq E$ are allowed (``independent''), with
\begin{itemize}
  \item $\emptyset$ is allowed;
  \item every subset of an allowed set is allowed;
  \item if $A$, $B$ are allowed and $|A| < |B|$, some $x \in B \setminus A$ makes $A \cup \{x\}$ allowed.
\end{itemize}
These are exactly the rules where greedy works: sort by weight, add an element if the set stays allowed.
\item Rules that are matroids: at most $cap_c$ elements of each colour $c$ (\texttt{ColorMat}); chosen edges form a forest (\texttt{GraphMat}); chosen vectors are linearly independent (xor basis); the unchosen edges keep the graph connected.
\item Matroid Intersection: given two such rules on the same $E$, it finds the largest $S$ allowed by both. Each rule is a struct with \texttt{build(S)} (called once per round) and \texttt{ind(x, y)} (is $S - x + y$ allowed? $x = -1$: $S + y$); xor basis and connectivity have no ready-made struct, write \texttt{ind} yourself. Typical uses:
\begin{itemize}
  \item spanning tree with at most $cap_c$ edges of colour $c$: forest $\cap$ colour caps; one exists iff the answer has $n - 1$ edges.
  \item one vector from each group, all independent: xor basis $\cap$ colour caps (group $=$ colour, $cap = 1$).
\end{itemize}
\end{itemize}
